\chapter{Diagramas detallados de arquitectura}\label{anexo:A}

Este anexo detalla las fronteras de confianza y la distribuci\'on de puertos y
adaptadores de la implementaci\'on. Las figuras complementan las vistas del
Cap\'itulo 3 y separan el equipo del titular, la aplicaci\'on y sus servicios.

\section{Vista detallada de Fronteras de Confianza}\label{sec:fronteras-confianza}

El navegador accede a Qadra mediante HTTPS. La API autentica con sesiones
opacas revocables almacenadas en Redis y aplica permisos por rol y expediente.
La aplicaci\'on conserva las operaciones y su auditor\'ia en PostgreSQL mediante
los adaptadores \texttt{postgres}/\texttt{tokio-postgres}; los archivos se
protegen con cifrado autenticado. La clave privada que el titular selecciona
para firmar localmente permanece en su equipo y se procesa en un worker
Rust/WASM. La CA tiene portal, proceso y acceso administrativo separados.
La TSA local y el proveedor de correo cumplen funciones acotadas; ninguno
convierte la confianza interna del prototipo en acreditaci\'on de identidad civil.

\begin{figure}[H]
  \centering
  \begin{tikzpicture}[
  box/.style={draw=ipnvino, rounded corners=1.5mm, align=center,
    text width=13cm, inner sep=7pt, font=\small},
  flow/.style={-{Stealth[length=2mm]}, line width=0.7pt}
]
  \node[box, fill=ipnvino!6] (browser) at (0,0) {
    \textbf{Equipo del titular: navegador Qadra}\\
    Clave cifrada y frase seleccionadas localmente; worker Rust/WASM.\\
    A Qadra se env\'ian certificados y firmas; el CSR se entrega a la CA.\\
    La clave privada y su frase permanecen locales.
  };
  \node[box, below=9mm of browser] (http) {
    \textbf{Frontera p\'ublica: HTTPS y API Axum}\\
    Validaci\'on de entradas, l\'imites y autenticaci\'on.\\
    Contrase\~na o certificado Owner habilitado, seguidos de MFA.\\
    Sesi\'on opaca revocable; permisos por rol y expediente.
  };
  \node[box, below=9mm of http] (app) {
    \textbf{Aplicaci\'on y dominio}\\
    Casos de uso, comprobaci\'on criptogr\'afica, historial y auditor\'ia.\\
    Autorizaci\'on y persistencia transaccional de la operaci\'on.
  };
  \node[box, below=9mm of app] (data) {
    \textbf{Persistencia privada y secretos del servidor}\\
    PostgreSQL: \texttt{postgres}/\texttt{tokio-postgres}, consultas parametrizadas.\\
    Redis: sesiones y estado temporal; archivos: cifrado AES-256-GCM.\\
    Claves de servicio y confianza configurada fuera del paquete de despliegue.
  };
  \node[box, dashed, below=9mm of data] (services) {
    \textbf{Servicios con autoridad y controles propios}\\
    CA: portal y proceso separados, acceso exclusivo de sus responsables.\\
    TSA local configurada: RFC 3161. Resend: entrega de correo.\\
    La CA interna y la TSA local no acreditan identidad civil ni un PSC externo.
  };
  \draw[flow] (browser) -- (http);
  \draw[flow] (http) -- (app);
  \draw[flow] (app) -- (data);
\end{tikzpicture}

  \caption{Fronteras de confianza, sesiones opacas y persistencia de la implementaci\'on.}
  \label{fig:fronteras-confianza}
\end{figure}

\section{Arquitectura hexagonal con puertos y adaptadores}\label{sec:hexagonal-completa}

Los transportes HTTP y CLI invocan los casos de uso; los binarios componen sus
dependencias concretas. Los traits de persistencia y de flujo residen en
\texttt{application}, mientras que \texttt{domain} define reglas y puertos
criptogr\'aficos y de auditor\'ia. \texttt{infrastructure} implementa los
adaptadores contra esas interfaces. La figura muestra ejemplos representativos,
sin equiparar cada handler de entrada con un trait ni suponer un transporte
WebSocket. Las flechas continuas muestran invocaci\'on o dependencia hacia el
interior; las discontinuas indican implementaci\'on de puertos.

\begin{figure}[H]
  \centering
  \begin{tikzpicture}[
  box/.style={draw=ipnvino, rounded corners=1.5mm, align=center,
    text width=13cm, inner sep=7pt, font=\small},
  flow/.style={-{Stealth[length=2mm]}, line width=0.7pt}
]
  \node[box] (entry) at (0,0) {
    \textbf{Entradas y composici\'on}\\
    HTTP: Axum en \texttt{web}; CLI: Clap en \texttt{bin}.\\
    \texttt{despacho-cli} y \texttt{qadra-ca} conectan servicios y adaptadores.
  };
  \node[box, fill=ipnvino!6, below=9mm of entry] (app) {
    \textbf{application: casos de uso y puertos de flujo/persistencia}\\
    Ejemplos: \texttt{HearingWorkflow}, \texttt{HearingStore},\\
    \texttt{IndividualSigningStore}, \texttt{DocumentSigningVerifier}.\\
    Los adaptadores de entrada invocan estos casos de uso.
  };
  \node[box, fill=ipnvino!6, below=9mm of app] (domain) {
    \textbf{domain: reglas y puertos de dominio}\\
    Ejemplos: \texttt{DocumentSigner}, \texttt{TimestampService},\\
    \texttt{AuthenticatedCipher}, \texttt{PasswordHasher}, \texttt{AuditLog}.
  };
  \node[box, below=9mm of domain] (infra) {
    \textbf{infrastructure: adaptadores de los puertos internos}\\
    PostgreSQL y Redis; archivos y auditor\'ia;\\
    criptograf\'ia, OpenSSL/PKI, TSA configurada y correo.\\
    Depende de \texttt{application} y \texttt{domain}; no a la inversa.
  };
  \draw[flow] (entry) -- (app);
  \draw[flow] (app) -- (domain);
  \draw[flow, dashed] (infra) -- (domain);
  \draw[flow, dashed] (infra.east) -- ++(0.5,0) |- (app.east);
\end{tikzpicture}

  \caption{Distribuci\'on de puertos entre application y domain, con adaptadores externos.}
  \label{fig:hexagonal-completa}
\end{figure}
